

\subsection{Transfer to video classification}
\paragraph{Experimental setup.}
We evaluate the quality of the representations learned by \chonk by adapting the model pretrained on images for video classification.
We follow the ``factorised encoder'' architecture of~\citet{arnab2021vivit}: Our video model consists of an initial ``spatial transformer'', which encodes each frame of the video independently of each other.
Thereafter, the representation from each frame is pooled into a single token, which is then fed to a subsequent ``temporal transformer'' that models the temporal relations between the representations of each frame.

% Our choice for this architecture is motivated by the results of~\citet{arnab2021vivit}, which showed that it was an effective method for adapting image-pretrained models for video classification tasks.
% The parameters of the ``spatial transformer'' are initialized from the model pretrained on images, while the temporal transformer is trained from random initialization. Furthermore, we 
Here, we initialize the ``spatial transformer'' with the pretrained weights from ViT-22B and freeze them, as this represents a computationally efficient method of adapting large-scale models for video, and also because it allows us to effectively evaluate the representations learned by pretraining \chonk. Exhaustive experimental details are included in \cref{sec:app_video_classification}. The temporal transformer is lightweight both in terms of parameters (only 63.7M parameters compared to the 22B frozen parameters in the spatial transformer), and FLOPs as it operates on a \emph{single} token per frame.

% \todo{Exhaustive experimental details are included in \cref{sec:app_video_classification}. The temporal transformer that we learn is comparatively lightweight, consisting of 63.7 million parameters, compared to the 22 billion frozen parameters in the frozen spatial transformer.}
\paragraph{Results.}
\begin{table}[t]
    \caption{
        Video classification results.
        We evaluate the \chonk representations by freezing the backbone, and training a small transformer to aggregate frozen, per-frame representations.
        \chonk outperforms the largest previous vision backbone, ViT-e~\citep{chen2022pali} which contains 4 billion parameters and is also pretrained on JFT.
    }
    \centering{
    \resizebox{.55\textwidth}{!}{%
    \begin{tabular}{@{}lcc@{}}
        \toprule
                             & Kinetics 400 & Moments in Time \\
        \midrule
        \multicolumn{3}{l}{\textit{Frozen backbone}}                   \\
        CoCA$^*$                 &   88.0           &    47.4             \\
        ViT-e               &   86.5           &     43.6           \\  % https://xm2a.corp.google.com/experiments/49222795/workUnits/8, http://xid/51624094/3
        \chonk              &   88.0           &    44.9           \\  % https://xm2a.corp.google.com/experiments/49276809/workUnits/5, https://xm2a.corp.google.com/experiments/49597259/workUnits/3
        \midrule
        Fully finetuned SOTA &   91.1           &    49.0              \\ % https://arxiv.org/pdf/2212.03191v2.pdf, https://arxiv.org/pdf/2205.01917v2.pdf
        \bottomrule
    \end{tabular}
    }}
    \begin{adjustwidth}{110pt}{110pt}
    \begin{spacing}{0.125}{\fontsize{8pt}{5pt}\selectfont $^*$Note that CoCA uses pre-pool spatial features and higher spatial resolution for both datasets. More details in \cref{sec:app_video_classification}.}\end{spacing}
    \end{adjustwidth}
    \label{tab:results_video}
    % \vspace{-10pt}
\end{table}
\cref{tab:results_video} presents our results on video classification on the Kinetics 400~\citep{kay2017kinetics} and Moments in Time~\citep{monfort2019moments} datasets, showing that we can achieve competitive results with a frozen backbone.
We first compare to ViT-e~\citep{chen2022pali}, which has the largest previous vision backbone model consisting of 4 billion parameters, and was also trained on the JFT dataset.
We observe that our larger \chonk model improves by 1.5 points on Kinetics 400, and 1.3 points on Moments in Time. 
Our results with a frozen backbone are also competitive with CoCA~\citep{yu2022coca}, which performs a combination of contrastive and generative caption pretraining in comparison to our supervised pretraining, and uses many tokens per frame (vs.\ a single one produced by the pretrained frozen pooling) as well as a higher testing resolution.

Finally, we note that there is headroom for further improvement by full end-to-end fine-tuning. This is evidenced by the current state-of-the-art on Kinetics 400~\citep{wang2022internvideo} and Moments in Time~\citep{yu2022coca} which leverage a combination of large-scale video pretraining and full end-to-end fine-tuning on the target dataset. 
%We leave these directions for future work.


